\documentclass[11pt]{article}
\usepackage[margin=25mm]{geometry}
\usepackage{fontspec}
\setmainfont{Latin Modern Roman}
\usepackage{amsmath,amssymb,mathtools}
\usepackage{unicode-math}
\setmathfont{Latin Modern Math}
\usepackage{xurl,hyperref}
\hypersetup{colorlinks=true,urlcolor=blue}
\setcounter{secnumdepth}{0}
\setlength{\parindent}{0pt}
\setlength{\parskip}{5pt}
\setlength{\emergencystretch}{3em}
\newcommand{\sref}[2]{\hyperlink{source:#1:#2}{[S#2]}}
\newcommand{\cataloguescope}[1]{\paragraph{Recorded target.}#1}
\newcommand{\R}{\mathbb R}
\newcommand{\C}{\mathbb C}
\newcommand{\N}{\mathbb N}
\newcommand{\Z}{\mathbb Z}
\newcommand{\Q}{\mathbb Q}
\newcommand{\F}{\mathbb F}
\newcommand{\eps}{\varepsilon}
\begin{document}
\section{Future Restrictions of Inner Automorphic Dilations}\label{statement}

\subsection{Definitions and mathematical statement}

\paragraph{Shared notation.}$\mathbb N=\{1,2,\ldots\}$; the standard number fields have their usual meanings. All additional conventions are stated locally.


\paragraph{Definitions and conventions.} A Markov semigroup $T_t$ is a continuous semigroup of unital completely positive maps on the specified operator algebra $B$. A dilation consists of an embedding $i:B\to A$, a conditional expectation or compression $p:A\to B$, and endomorphisms $\vartheta_t$ with $T_t=p\vartheta_t i$. The GNS construction for $p$ gives a Hilbert $B$-module $E$ and a represented algebra in $\mathcal B^a(E)$. An automorphic dilation has invertible $*$-automorphisms; an inner one is implemented by unitaries. With cyclic reference vector $\xi$, the future module is the closed span $E_+=\overline{\operatorname{span}}\{\vartheta_t(i(b))\xi:t\ge0,\ b\in B\}$ in the source's construction. A proper $E_0$-semigroup here is the resulting genuinely endomorphic future dynamics rather than the original two-sided automorphic flow. A unital embedding of $\mathcal B^a(E_+)$ into $\mathcal B^a(E)$ must be compatible with the given flow, not merely an abstract inclusion.

\paragraph{Source scope and qualifications.} Retain the restriction, extension, and unital compatible embedding as separate requirements, including the proposed two-sided q-Fock case study. The source notes that the embedding would correspond to a module factorization E=E\_+ tensor F; it is not assumed.

\paragraph{Catalogue mathematical target.}
\paragraph{Catalogue wording retained for traceability.}
Under suitable minimality conditions, can every inner automorphic dilation be restricted to a proper $E_0$-semigroup on the future part $E_+$ of $\mathcal B^a(E)$? Does it extend to all of $\mathcal B^a(E_+)$, and can $\mathcal B^a(E_+)$ be unitally embedded in $\mathcal B^a(E)$ so that the extension is the restriction of the original automorphism semigroup? The two-sided time shift on $q$-Fock space is proposed as a case study.

\paragraph{Saved literature context.} The supplied catalogue's assessment dated 2026-08-21 records: The proposed restriction and factorization for inner automorphic dilations, including the two-sided q-Fock case, remains open; later dilation results do not supply the requested embedding and extension. The assessment is retained as dated provenance; it is not a new comprehensive literature or open-status certification.
\paragraph{Scope and source qualification.}Retain the restriction, extension, and unital compatible embedding as separate requirements, including the proposed two-sided q-Fock case study. The source notes that the embedding would correspond to a module factorization E=E_+ tensor F; it is not assumed.
\subsection{Short English statement}
Restrict a minimal inner dilation to its future module, extend that future dynamics, and embed it compatibly into the original dilation.
\subsection{Research attempt}
\paragraph{Model, date and investigated scope.}
OpenAI GPT-6 Astra (configured model \texttt{gpt-6-astra}), ultra reasoning; generated 11 October 2026. Proves explicit compression and dimension-divisibility obstructions to the naive future-subspace construction.

\paragraph{Sources and literature check.}
The complete supplied definition, including source qualifications and the dated literature assessment, was reviewed. The calculation below is an elementary restricted attempt motivated by [S2]; it does not claim novelty or a fresh literature-status determination.

\paragraph{Context and investigated scope.} The complete supplied definition, including source qualifications and the dated literature assessment, was reviewed. The calculation below is an elementary restricted attempt motivated by [S2]; it does not claim novelty or a fresh literature-status determination.

\paragraph{Compression can destroy multiplicativity.} Let $\mathcal H_+\subset\mathcal H$ be a proper closed subspace with orthogonal projection $P$. Even if a unitary flow is defined on $\mathcal H$, compression $X\mapsto PXP|_{\mathcal H_+}$ is generally not a homomorphism. On $\mathbb C^2$, take $P=|e_1\rangle\langle e_1|$ and $X=\left(\begin{smallmatrix}0&1\\1&0\end{smallmatrix}\right)$. Then $PXP=0$ but $PX^2P=P\ne0$. Thus a future restriction cannot be obtained by arbitrary compression of the original automorphism algebra.

\paragraph{A finite-dimensional embedding obstruction.} If a unital $*$-embedding $M_m(\mathbb C)\hookrightarrow M_n(\mathbb C)$ exists, the images of the $m$ diagonal matrix units are orthogonal projections summing to $I_n$. Off-diagonal matrix units give partial isometries between their ranges, so all ranges have the same dimension $k$. Hence $n=mk$. For example a two-dimensional future subspace inside $\mathbb C^3$ cannot yield a unital embedding of its full operator algebra into $M_3$, despite the obvious Hilbert-space inclusion.

\paragraph{Where the source problem starts.} These examples are not minimal inner Markov dilations satisfying all the source assumptions. They show that future-subspace inclusion alone supplies neither an endomorphism nor the required unital compatible embedding; a factorization or equivalent representation structure is essential. Establishing it under meaningful minimality, extending the future dynamics, and analyzing the suggested two-sided $q$-Fock shift remain separate gaps. No $q$-Fock factorization is inferred from its tensor-Fock analogue.

\paragraph{Verification and adversarial audit.}
Checked the 2×2 compression calculation and matrix-unit divisibility proof; did not assert that the toy dimensions form a source dilation. Model self-audit only; no independent human proof checking. Prompts A and B were applied to this canonical record and its qualified source scope. Checked the 2×2 compression calculation and matrix-unit divisibility proof; did not assert that the toy dimensions form a source dilation. Retained unresolved status and zero original-target progress; The required compatible module factorization/unital embedding and the q-Fock case have not been established under source minimality.

Coordinator checked the compression counterexample and matrix-unit rank argument for the unital embedding obstruction. This additional check was model-assisted, not independent human verification.

\paragraph{Tools, process and contributions.}
Python 3 was used for file preparation and structural checks; no numerical experiment or theorem prover was used. Contribution prompt A was followed by prompt B. GPT-6 Astra with ultra reasoning drafted the research and performed its stated self-audit; any additional coordinator contribution is identified in the verification and audit. A GPT-6-based Codex coordinator assembled the manuscript and checked repository integration; its exact variant and reasoning configuration were not exposed. No human mathematical contribution, independent human proof check or formal proof certification is supplied. Complete derivations and the consulted source references are included here; private reasoning transcripts are not published.

\paragraph{Final status and first remaining gap.}
\textbf{UNRESOLVED}. The required compatible module factorization/unital embedding and the q-Fock case have not been established under source minimality.

\section{Completion Estimate}
\noindent\textbf{COMPLETION: 0\%}
The calculation is an auxiliary or restricted attempt and establishes no newly verified portion of the full source target. The required compatible module factorization/unital embedding and the q-Fock case have not been established under source minimality.

\subsection{Sources}
\begin{itemize}
\item[\hypertarget{source:statement:1}{S1}] ULAM.AI and the UnsolvedMath Contributors, UnsolvedMath, catalogue entry. Catalogue curation under CC BY 4.0; definitions and English exposition added locally.
\url{https://huggingface.co/datasets/ulamai/UnsolvedMath}.
\item[\hypertarget{source:statement:2}{S2}] Mini-Workshop: Product Systems and Independence in Quantum Dynamics (2009), Oberwolfach Reports, DOI 10.4171/owr/2009/09, pp. 493--548; official report source and mathematical context.
\url{https://doi.org/10.4171/owr/2009/09}.
\item[\hypertarget{source:statement:3}{S3}] Official EMS Press PDF of the same Oberwolfach report; PDF page 51, printed page 543.
\url{https://ems.press/content/serial-article-files/46209}.
\item[\hypertarget{source:statement:4}{S4}] spatial Markov semigroups and dilations; reference in the saved dated catalogue literature review.
\url{https://doi.org/10.3842/SIGMA.2022.071}.
\end{itemize}
\label{end:statement}
\end{document}
